\answer{ex:20:1}{$t_{\text{v}}=\tau_r\ln\frac{1-L}{1-H}=16.8$~h, $t_{\text{s}}=\tau_d\ln\frac HL=5.8$~h, période $22.5$~h; c'est l'oscillation circadienne des seuils qui amène l'oscillateur à $24$~h: un verrouillage de phase~\eqref{eq:pll}.}
\answer{ex:20:2}{(a)~$L=0.111$; avec $L=0.092$, le sommeil durerait $9.6$~h. (b)~De $22\%$ ($1/0.82=1.22$), et non de $18\%$.}
\answer{ex:20:3}{$H=\frac{p-1}{p-1/q}=0.623$, $L=H/q=0.093$, où $p=e^{16/\tau_r}$, $q=e^{8/\tau_d}$. Après un réveil au bout de $4$~h, la phase est décalée pour toujours (endormissement $7.2$~h plus tôt); avec des seuils oscillants, elle revient en deux ou trois jours.}
\answer{ex:20:4}{(a)~$r_\pm=\frac12\bigl(1\pm\sqrt{1-4k/w}\bigr)=0.947,\ 0.053$; $a_{\text{h}}=3.51$, $a_{\text{b}}=0.49$. (b)~Activité $\approx0.33$~s, silence $\approx0.59$~s, période $\approx0.93$~s, fraction d'activité $0.36$.}
\answer{ex:20:5}{$a_{\text{h}}/r_+<b<a_{\text{b}}/r_-$: $3.71<b<9.20$. \nt{ACET} fait descendre $b$ sous $3.71$: l'intersection passe sur la branche haute, et c'est l'état actif ($r\approx1$) qui est stable.}
\answer{ex:20:6}{$4.9$, $6.8$, $8.8$, $5.5$, $8.9$~h pour $D=24$, $32$, $40$, $48$, $64$: la durée est fixée par la phase circadienne de l'endormissement, et non par la dette.}
\answer{ex:20:7}{Après B, l'erreur sur A vaut en moyenne $0.51$ (de $0.19$ à $1.08$ sur $20$ tirages; une réponse nulle donnerait $1$). Avec entrelacement, les deux erreurs sont inférieures à $10^{-4}$.}
\answer{ex:20:8}{Question ouverte. Une mise à l'échelle uniforme conserve les rapports des poids, mais ne conserve les réponses d'un neurone à seuil que si le seuil est mis à l'échelle du même facteur; les répétitions entrelacées modifient les rapports. La stabilité des gros contacts contredit une mise à l'échelle purement uniforme.}
